\documentclass[UTF8,12pt]{ctexart}
\usepackage[a4paper,margin=2.3cm]{geometry}
\usepackage{amsmath,amssymb,booktabs,graphicx,float,array}
\setlength{\parindent}{2em}
\title{基于季节自适应预测与分位数风险控制的微网日前购电模型}
\author{}
\date{}

\begin{document}
\maketitle

\begin{abstract}
针对负载与光伏逐日变化、日前计划一次确定且缺口电量按实时电价5倍采购的问题，本文建立“因果预测—风险修正—储能调度”模型。预测层采用近期历史、岭回归和浅层LightGBM在线集成，并以历史光伏有效窗口和日电量比实施季节水平校正；残余星期效应不显著且会增加回测成本，故最终关闭周效应。风险层根据过去90天净负荷残差构造逐时段分位数安全余量，分位数取$q=0.80$，该值直接来自5倍紧急电价下的报童临界比$4c/(4c+c)$，为事前给定的理论值、不由评价期回测挑选；决策层建立计划购电、充放电、弃光和储电状态的线性规划。2025年2月1日至12月31日334天回放中，净负荷MAE为247.43 kW，全年总费用为13918258.75元，紧急购电80111.20 kWh。全部48096个评价时段均满足物理约束。

\textbf{关键词：}微网调度；分位数风险；储能；线性规划；因果预测
\end{abstract}

\section{问题重述与分析}
微网每天0:00确定144个十分钟时段的计划购电量。计划电量无论是否使用均按原价结算；执行时若计划购电、光伏和储能放电仍不能满足负载，则缺口按对应电价5倍紧急采购。因此预测偏低与偏高具有显著不对称代价，且储能使相邻时段的决策相互关联。

本文采用三层结构：首先仅用预测日前历史数据估计负载和光伏；其次由历史净负荷残差给预测附加风险余量；最后将修正后的净需求输入储能线性规划。取$E_{1,0}=E_0=6000$ kWh，并按$E_{d+1,0}=E_{d,T}$逐日连续递推；1月用于预测预热和状态递推，2月1日至12月31日为正式费用评价期。

\section{假设与符号}
假设十分钟时段内功率和电价不变；第$d$天决策只使用$j<d$的历史；充放电效率均为0.9，储能范围为1200--10800 kWh，功率上限为5000 kW；题目未提供的气象和位置变量不进入模型。

记$L_{d,t},P_{d,t}^{\rm act}$为负载和光伏实际功率，$L_{d,t}^{\rm fc},P_{d,t}^{\rm fc}$为日前预测功率，$m_{d,t}$为净负荷安全余量；$x,y,r$分别为计划购电、实际使用计划购电和购电浪费功率，$g,w$为光伏实际使用和光伏浪费功率，$u,v$为充、放电功率，$e$为紧急购电功率，$E$为储电量，$c_t$为问题2的固定日内电价曲线，$\Delta t=1/6$小时。所有决策流量均以kW表示，结果表中的时段电量为功率乘$\Delta t$。

\section{因果预测模型}
对负载与光伏分别构造近期加权基线、岭回归和浅层LightGBM，使用1、2、7、14日滞后、近3/7/14日同一时段均值、日内周期和近期水平趋势。候选模型按近期指数加权误差组合，训练窗随历史残差漂移在45--180天内调整。

根据历史光伏曲线估计有效发电窗口$W_d$，窗口外预测置零；再由历史实际日电量与基础预测日电量之比估计季节水平因子$\lambda_d$：
\[
\lambda_d=\operatorname{clip}\left(1+\frac{n}{n+\kappa}
\left[\frac{\sum_j\omega_j\rho_j}{\sum_j\omega_j}-1\right],0.85,1.15\right).
\]
最终有
\[
L_{d,t}^{\rm fc}=\lambda_d^L L_{d,t}^{\rm fc,(0)},\qquad
P_{d,t}^{\rm fc}=\lambda_d^P P_{d,t}^{\rm fc,(0)}\mathbf1(t\in W_d).
\]
星期效应检验$p=0.967$；加入周效应后净负荷MAE从247.43升至247.84 kW、费用增加14992.74元，故最终模型删除该通道。

\begin{table}[H]\centering
\caption{正式评价期预测精度}
\begin{tabular}{lrrrr}\toprule
目标 & MAE/kW & RMSE/kW & nMAE & 偏差/kW\\\midrule
负载 & 148.70 & 197.72 & 3.22\% & 2.10\\
光伏 & 149.88 & 304.11 & 6.30\% & 3.09\\
净负荷 & 247.43 & 362.77 & 8.25\% & -0.99\\\bottomrule
\end{tabular}
\end{table}

\section{风险控制与日前调度}
令$N_{d,t}^{\rm fc}=L_{d,t}^{\rm fc}-P_{d,t}^{\rm fc}$，净负荷误差为$\xi_{j,t}=(L_{j,t}-P_{j,t}^{\rm act})-N_{j,t}^{\rm fc}$。第$d$天只用$j<d$的误差，计算90天长窗与30天短窗的$q$分位数：
\[
\widetilde m_{d,t}=0.7Q_q(\xi_{d-90:d-1,t})+0.3Q_q(\xi_{d-30:d-1,t}),\qquad
m_{d,t}=0.35\widetilde m_{d,t}+0.65m_{d-1,t}.
\]
分位数取$q=0.80$，直接来自5倍紧急电价下的报童临界比$4c/(4c+c)$，不含待估参数；后文的风险分位数灵敏度网格仅作事后对照，不参与选参。计划净需求为$N_{d,t}^{\rm plan}=N_{d,t}^{\rm fc}+m_{d,t}$。

每日求解线性规划
\[
\min \sum_{t=1}^{T}c_tx_{d,t}\Delta t-\lambda_EE_{d,T}
+\varepsilon_0\sum_t(u_{d,t}+v_{d,t}+w_{d,t})\Delta t,
\]
满足
\[
x_{d,t}+P_{d,t}^{\rm fc}-w_{d,t}+v_{d,t}
=L_{d,t}^{\rm fc}+m_{d,t}+u_{d,t},
\]
\[
E_{d,t}=E_{d,t-1}+0.9u_{d,t}\Delta t-v_{d,t}\Delta t/0.9,
\]
以及储能容量、充放电功率和非负约束。日末水值$\lambda_E=0.9\bar c$用于连接相邻日期。实际执行满足
\[
y_{d,t}+g_{d,t}+v_{d,t}+e_{d,t}=L_{d,t}+u_{d,t},\quad
g_{d,t}+w_{d,t}=P_{d,t}^{\rm act},\quad
r_{d,t}=x_{d,t}-y_{d,t}\ge0,
\]
其中$e_{d,t}$为紧急购电功率。总费用为
\[
C=\sum_{d=32}^{D}\sum_{t=1}^{T}(c_tx_{d,t}+5c_te_{d,t})\Delta t.
\]
其中$C_d^{\rm grid}=\sum_t c_tx_{d,t}\Delta t$，$C_d^{\rm emer}=\sum_t5c_te_{d,t}\Delta t$，且$C_d=C_d^{\rm grid}+C_d^{\rm emer}$。

\section{全年结果}
\begin{table}[H]\centering
\caption{2025年2月1日至12月31日结果}
\begin{tabular}{lr}\toprule
指标 & 数值\\\midrule
计划购电量/kWh & 22011526.88\\
计划购电费/元 & 13418699.15\\
紧急购电量/kWh & 80111.20\\
紧急购电费/元 & 499559.60\\
总费用/元 & \textbf{13918258.75}\\
未用计划电量/kWh & 1467144.65\\
弃光量/kWh & 1401730.20\\
紧急购电日期数 & 124\\
单日费用CVaR95/元 & 71000.05\\\bottomrule
\end{tabular}
\end{table}

\section{四个指定日期结果}
\begin{table}[H]\centering\small
\caption{表1格式的计划购电结果（电量kWh，费用元）}
\resizebox{\textwidth}{!}{\begin{tabular}{lrrrrrrrr}\toprule
日期&10:00&12:00&14:00&16:00&18:00&20:00&全天电量&全天费用\\\midrule
3月20日&0.00&625.33&0.00&494.01&720.26&0.00&69732.29&42273.50\\
6月21日&0.00&0.00&0.00&238.56&0.00&0.00&37971.96&21833.00\\
9月23日&0.00&534.14&0.00&605.30&813.43&0.00&73786.83&46091.89\\
12月21日&0.00&974.72&0.00&830.86&724.94&0.00&97558.19&62541.25\\\bottomrule
\end{tabular}}
\end{table}

\begin{table}[H]\centering\small
\caption{表2格式的储能结果（各区间为充/放，单位kWh）}
\resizebox{\textwidth}{!}{\begin{tabular}{lrrrrrrrr}\toprule
日期&0--4&4--8&8--12&12--16&16--20&20--24&0:00储电&24:00储电\\\midrule
3月20日&88.51/108.15&165.81/5434.43&6285.65/2777.38&4046.04/450.44&320.03/5534.25&10371.56/2930.22&10800&10800\\
6月21日&75.26/3649.47&379.11/4836.61&10022.28/0&0/0&170.75/5457.79&9636.27/2485.90&10800&10800\\
9月23日&0.28/0.23&312.19/5274.70&6212.68/1896.62&2885.44/451.03&0/5532.04&10354.10/2854.78&10800&10800\\
12月21日&55.38/134.20&953.21/1503.81&5562.43/7241.30&7136.25/2226.67&24.51/5743.00&10578.73/2842.53&10800&10800\\\bottomrule
\end{tabular}}
\end{table}

\begin{table}[H]\centering
\caption{指定日期紧急购电量}
\begin{tabular}{lrr}\toprule
日期&紧急购电/kWh&紧急购电费/元\\\midrule
2025-03-20&0.000&0.00\\
2025-06-21&0.000&0.00\\
2025-09-23&1.888&12.73\\
2025-12-21&0.000&0.00\\\bottomrule
\end{tabular}
\end{table}

\begin{figure}[H]\centering
\includegraphics[width=0.95\textwidth]{../figures/指定日期调度.png}
\caption{四个指定日期的负载、光伏及购电策略}
\end{figure}

\section{灵敏度与稳健性}
\begin{table}[H]\centering
\caption{风险分位数灵敏度}
\begin{tabular}{rrrrr}\toprule
$q$&开发期费用/元&留出期费用/元&全年费用/元&紧急日数\\\midrule
0.60&8917123.60&5344179.65&14261303.25&255\\
0.65&8730343.75&5287893.88&14018237.63&221\\
0.70&8620409.52&\textbf{5278910.66}&13899320.18&188\\
0.75&\textbf{8568601.01}&5295709.78&13864310.79&150\\
\textbf{0.80}（采用）&8585080.60&5333178.16&\textbf{13918258.75}&124\\
0.85&8652894.21&5400908.26&14053802.47&92\\\bottomrule
\end{tabular}
\end{table}

本模型采用理论临界比0.80。作为事后对照，开发期最优$q=0.75$、留出期最优点为0.70，说明最优点位置随划分区间漂移；从0.75提高到0.80使费用增加53947.96元（+0.39%）、紧急日减少26天，对应每减少一个紧急日约2075元。该网格不参与选参：若改用开发期挑出的0.75，就等于用评价窗口内部的数据反推风险参数。

全年储能越界、功率越界、同时充放电和负流量次数均为0，最大能量平衡误差为$2.27\times10^{-13}$ kWh。未来扰动前缀复验中，负载和光伏预测最大变化均为0。

\section{评价与结论}
模型严格遵守日前因果信息集，只保留有样本外收益的季节通道；线性规划规模小、可解释且易复现；残差分位数随预测质量自动收缩，能够将预测改进传导到决策。局限是数据仅覆盖一年且无天气预报，对极端光伏误差的解释能力有限；全局$q$也尚未按季节、电价与SOC分别设置。

最终采用“光伏窗口与季节水平校正＋0.80因果分位数安全余量＋确定性储能LP”。该模型全年总费用为13918258.75元，四个指定日期及全年逐十分钟策略均由同一程序生成。

\end{document}
